% Ngân hàng câu hỏi TOÁN THỰC TẾ T10 -- Chương 6, Bài 17, Dạng D73
%% Dạng: Bất phương trình bậc hai trong bài toán kinh doanh
%% 100 câu: NB 30 -- TH 40 -- VD 20 -- VDC 10 ; mỗi mức: trắc nghiệm : đúng/sai : trả lời ngắn = 3 : 1 : 1
%% Phương án ĐẦU TIÊN trong \cauTN luôn là đáp án đúng (gói trondethi tự đảo).
%% Phần II: \y{nội dung}{d|s}. Phần III: đáp án đặt ở đối số cuối.

\cauTN[id=T10-C6-B17-D73-NB-001,muc=NB,loigiai={Doanh nghiệp không lỗ khi $L(x)\ge0 \Leftrightarrow x^2-33x+200\le0$. Tam thức có hai nghiệm $x=8$ và $x=25$, hệ số $a>0$ nên nghiệm của bất phương trình là $8\le x\le25$.}]{Một xưởng cơ khí có lợi nhuận khi nhận $x$ đơn hàng mỗi tháng là $L(x)=-x^2+33x-200$ (triệu đồng). Hỏi xưởng cần nhận bao nhiêu đơn hàng để không bị lỗ?}
  {$8\le x\le25$.}
  {$x\le8$ hoặc $x\ge25$.}{$8<x<27$.}{$x\ge25$.}

\cauTN[id=T10-C6-B17-D73-NB-002,muc=NB,loigiai={Doanh nghiệp không lỗ khi $L(x)\ge0 \Leftrightarrow x^2-69x+1178\le0$. Tam thức có hai nghiệm $x=31$ và $x=38$, hệ số $a>0$ nên nghiệm của bất phương trình là $31\le x\le38$.}]{Một xưởng cơ khí có lợi nhuận khi nhận $x$ đơn hàng mỗi tháng là $L(x)=-x^2+69x-1178$ (triệu đồng). Hỏi xưởng cần nhận bao nhiêu đơn hàng để không bị lỗ?}
  {$31\le x\le38$.}
  {$x\le31$ hoặc $x\ge38$.}{$31<x<40$.}{$x\ge38$.}

\cauTN[id=T10-C6-B17-D73-NB-003,muc=NB,loigiai={Doanh nghiệp không lỗ khi $L(x)\ge0 \Leftrightarrow x^2-46x+360\le0$. Tam thức có hai nghiệm $x=10$ và $x=36$, hệ số $a>0$ nên nghiệm của bất phương trình là $10\le x\le36$.}]{Một doanh nghiệp có lợi nhuận khi sản xuất $x$ sản phẩm là $L(x)=-x^2+46x-360$ (triệu đồng). Hỏi doanh nghiệp cần sản xuất bao nhiêu sản phẩm để không bị lỗ?}
  {$10\le x\le36$.}
  {$x\le10$ hoặc $x\ge36$.}{$10<x<38$.}{$x\ge36$.}

\cauTN[id=T10-C6-B17-D73-NB-004,muc=NB,loigiai={Doanh nghiệp không lỗ khi $L(x)\ge0 \Leftrightarrow 2x^2-78x+580\le0$. Tam thức có hai nghiệm $x=10$ và $x=29$, hệ số $a>0$ nên nghiệm của bất phương trình là $10\le x\le29$.}]{Một xưởng cơ khí có lợi nhuận khi nhận $x$ đơn hàng mỗi tháng là $L(x)=-2x^2+78x-580$ (triệu đồng). Hỏi xưởng cần nhận bao nhiêu đơn hàng để không bị lỗ?}
  {$10\le x\le29$.}
  {$x\le10$ hoặc $x\ge29$.}{$10<x<31$.}{$x\ge29$.}

\cauTN[id=T10-C6-B17-D73-NB-005,muc=NB,loigiai={Doanh nghiệp không lỗ khi $L(x)\ge0 \Leftrightarrow 2x^2-98x+1020\le0$. Tam thức có hai nghiệm $x=15$ và $x=34$, hệ số $a>0$ nên nghiệm của bất phương trình là $15\le x\le34$.}]{Một doanh nghiệp có lợi nhuận khi sản xuất $x$ sản phẩm là $L(x)=-2x^2+98x-1020$ (triệu đồng). Hỏi doanh nghiệp cần sản xuất bao nhiêu sản phẩm để không bị lỗ?}
  {$15\le x\le34$.}
  {$x\le15$ hoặc $x\ge34$.}{$15<x<36$.}{$x\ge34$.}

\cauTN[id=T10-C6-B17-D73-NB-006,muc=NB,loigiai={Doanh nghiệp không lỗ khi $L(x)\ge0 \Leftrightarrow 2x^2-60x+378\le0$. Tam thức có hai nghiệm $x=9$ và $x=21$, hệ số $a>0$ nên nghiệm của bất phương trình là $9\le x\le21$.}]{Một doanh nghiệp có lợi nhuận khi sản xuất $x$ sản phẩm là $L(x)=-2x^2+60x-378$ (triệu đồng). Hỏi doanh nghiệp cần sản xuất bao nhiêu sản phẩm để không bị lỗ?}
  {$9\le x\le21$.}
  {$x\le9$ hoặc $x\ge21$.}{$9<x<23$.}{$x\ge21$.}

\cauTN[id=T10-C6-B17-D73-NB-007,muc=NB,loigiai={Doanh nghiệp không lỗ khi $L(x)\ge0 \Leftrightarrow x^2-30x+81\le0$. Tam thức có hai nghiệm $x=3$ và $x=27$, hệ số $a>0$ nên nghiệm của bất phương trình là $3\le x\le27$.}]{Một xưởng cơ khí có lợi nhuận khi nhận $x$ đơn hàng mỗi tháng là $L(x)=-x^2+30x-81$ (triệu đồng). Hỏi xưởng cần nhận bao nhiêu đơn hàng để không bị lỗ?}
  {$3\le x\le27$.}
  {$x\le3$ hoặc $x\ge27$.}{$3<x<29$.}{$x\ge27$.}

\cauTN[id=T10-C6-B17-D73-NB-008,muc=NB,loigiai={Doanh nghiệp không lỗ khi $L(x)\ge0 \Leftrightarrow 2x^2-90x+648\le0$. Tam thức có hai nghiệm $x=9$ và $x=36$, hệ số $a>0$ nên nghiệm của bất phương trình là $9\le x\le36$.}]{Một doanh nghiệp có lợi nhuận khi sản xuất $x$ sản phẩm là $L(x)=-2x^2+90x-648$ (triệu đồng). Hỏi doanh nghiệp cần sản xuất bao nhiêu sản phẩm để không bị lỗ?}
  {$9\le x\le36$.}
  {$x\le9$ hoặc $x\ge36$.}{$9<x<38$.}{$x\ge36$.}

\cauTN[id=T10-C6-B17-D73-NB-009,muc=NB,loigiai={Doanh nghiệp không lỗ khi $L(x)\ge0 \Leftrightarrow 2x^2-50x+300\le0$. Tam thức có hai nghiệm $x=10$ và $x=15$, hệ số $a>0$ nên nghiệm của bất phương trình là $10\le x\le15$.}]{Một hợp tác xã trồng rau ở Củ Chi có lợi nhuận khi canh tác $x$ sào đất là $L(x)=-2x^2+50x-300$ (triệu đồng). Hỏi hợp tác xã nên canh tác bao nhiêu sào để không bị lỗ?}
  {$10\le x\le15$.}
  {$x\le10$ hoặc $x\ge15$.}{$10<x<17$.}{$x\ge15$.}

\cauTN[id=T10-C6-B17-D73-NB-010,muc=NB,loigiai={Doanh nghiệp không lỗ khi $L(x)\ge0 \Leftrightarrow 2x^2-100x+1200\le0$. Tam thức có hai nghiệm $x=20$ và $x=30$, hệ số $a>0$ nên nghiệm của bất phương trình là $20\le x\le30$.}]{Một hợp tác xã trồng rau ở Củ Chi có lợi nhuận khi canh tác $x$ sào đất là $L(x)=-2x^2+100x-1200$ (triệu đồng). Hỏi hợp tác xã nên canh tác bao nhiêu sào để không bị lỗ?}
  {$20\le x\le30$.}
  {$x\le20$ hoặc $x\ge30$.}{$20<x<32$.}{$x\ge30$.}

\cauTN[id=T10-C6-B17-D73-NB-011,muc=NB,loigiai={Doanh nghiệp không lỗ khi $L(x)\ge0 \Leftrightarrow x^2-43x+306\le0$. Tam thức có hai nghiệm $x=9$ và $x=34$, hệ số $a>0$ nên nghiệm của bất phương trình là $9\le x\le34$.}]{Một xưởng cơ khí có lợi nhuận khi nhận $x$ đơn hàng mỗi tháng là $L(x)=-x^2+43x-306$ (triệu đồng). Hỏi xưởng cần nhận bao nhiêu đơn hàng để không bị lỗ?}
  {$9\le x\le34$.}
  {$x\le9$ hoặc $x\ge34$.}{$9<x<36$.}{$x\ge34$.}

\cauTN[id=T10-C6-B17-D73-NB-012,muc=NB,loigiai={Doanh nghiệp không lỗ khi $L(x)\ge0 \Leftrightarrow x^2-40x+364\le0$. Tam thức có hai nghiệm $x=14$ và $x=26$, hệ số $a>0$ nên nghiệm của bất phương trình là $14\le x\le26$.}]{Một quán ăn có lợi nhuận trong ngày khi phục vụ $x$ lượt khách là $L(x)=-x^2+40x-364$ (nghìn đồng). Hỏi quán cần phục vụ bao nhiêu lượt khách để không bị lỗ?}
  {$14\le x\le26$.}
  {$x\le14$ hoặc $x\ge26$.}{$14<x<28$.}{$x\ge26$.}

\cauTN[id=T10-C6-B17-D73-NB-013,muc=NB,loigiai={Doanh nghiệp không lỗ khi $L(x)\ge0 \Leftrightarrow x^2-47x+420\le0$. Tam thức có hai nghiệm $x=12$ và $x=35$, hệ số $a>0$ nên nghiệm của bất phương trình là $12\le x\le35$.}]{Một hợp tác xã trồng rau ở Củ Chi có lợi nhuận khi canh tác $x$ sào đất là $L(x)=-x^2+47x-420$ (triệu đồng). Hỏi hợp tác xã nên canh tác bao nhiêu sào để không bị lỗ?}
  {$12\le x\le35$.}
  {$x\le12$ hoặc $x\ge35$.}{$12<x<37$.}{$x\ge35$.}

\cauTN[id=T10-C6-B17-D73-NB-014,muc=NB,loigiai={Doanh nghiệp không lỗ khi $L(x)\ge0 \Leftrightarrow 2x^2-76x+210\le0$. Tam thức có hai nghiệm $x=3$ và $x=35$, hệ số $a>0$ nên nghiệm của bất phương trình là $3\le x\le35$.}]{Một hợp tác xã trồng rau ở Củ Chi có lợi nhuận khi canh tác $x$ sào đất là $L(x)=-2x^2+76x-210$ (triệu đồng). Hỏi hợp tác xã nên canh tác bao nhiêu sào để không bị lỗ?}
  {$3\le x\le35$.}
  {$x\le3$ hoặc $x\ge35$.}{$3<x<37$.}{$x\ge35$.}

\cauTN[id=T10-C6-B17-D73-NB-015,muc=NB,loigiai={Doanh nghiệp không lỗ khi $L(x)\ge0 \Leftrightarrow x^2-60x+851\le0$. Tam thức có hai nghiệm $x=23$ và $x=37$, hệ số $a>0$ nên nghiệm của bất phương trình là $23\le x\le37$.}]{Một xưởng cơ khí có lợi nhuận khi nhận $x$ đơn hàng mỗi tháng là $L(x)=-x^2+60x-851$ (triệu đồng). Hỏi xưởng cần nhận bao nhiêu đơn hàng để không bị lỗ?}
  {$23\le x\le37$.}
  {$x\le23$ hoặc $x\ge37$.}{$23<x<39$.}{$x\ge37$.}

\cauTN[id=T10-C6-B17-D73-NB-016,muc=NB,loigiai={Doanh nghiệp không lỗ khi $L(x)\ge0 \Leftrightarrow 2x^2-80x+728\le0$. Tam thức có hai nghiệm $x=14$ và $x=26$, hệ số $a>0$ nên nghiệm của bất phương trình là $14\le x\le26$.}]{Một hợp tác xã trồng rau ở Củ Chi có lợi nhuận khi canh tác $x$ sào đất là $L(x)=-2x^2+80x-728$ (triệu đồng). Hỏi hợp tác xã nên canh tác bao nhiêu sào để không bị lỗ?}
  {$14\le x\le26$.}
  {$x\le14$ hoặc $x\ge26$.}{$14<x<28$.}{$x\ge26$.}

\cauTN[id=T10-C6-B17-D73-NB-017,muc=NB,loigiai={Doanh nghiệp không lỗ khi $L(x)\ge0 \Leftrightarrow 2x^2-116x+1632\le0$. Tam thức có hai nghiệm $x=24$ và $x=34$, hệ số $a>0$ nên nghiệm của bất phương trình là $24\le x\le34$.}]{Một hợp tác xã trồng rau ở Củ Chi có lợi nhuận khi canh tác $x$ sào đất là $L(x)=-2x^2+116x-1632$ (triệu đồng). Hỏi hợp tác xã nên canh tác bao nhiêu sào để không bị lỗ?}
  {$24\le x\le34$.}
  {$x\le24$ hoặc $x\ge34$.}{$24<x<36$.}{$x\ge34$.}

\cauTN[id=T10-C6-B17-D73-NB-018,muc=NB,loigiai={Doanh nghiệp không lỗ khi $L(x)\ge0 \Leftrightarrow 2x^2-48x+280\le0$. Tam thức có hai nghiệm $x=10$ và $x=14$, hệ số $a>0$ nên nghiệm của bất phương trình là $10\le x\le14$.}]{Một hợp tác xã trồng rau ở Củ Chi có lợi nhuận khi canh tác $x$ sào đất là $L(x)=-2x^2+48x-280$ (triệu đồng). Hỏi hợp tác xã nên canh tác bao nhiêu sào để không bị lỗ?}
  {$10\le x\le14$.}
  {$x\le10$ hoặc $x\ge14$.}{$10<x<16$.}{$x\ge14$.}

\cauDS[id=T10-C6-B17-D73-NB-019,muc=NB,loigiai={Doanh nghiệp không lỗ khi $L(x)\ge0 \Leftrightarrow 2x^2-38x+168\le0$. Tam thức có hai nghiệm $x=7$ và $x=12$, hệ số $a>0$ nên nghiệm của bất phương trình là $7\le x\le12$.}]{Một doanh nghiệp có lợi nhuận khi sản xuất $x$ sản phẩm là $L(x)=-2x^2+38x-168$ (triệu đồng).}{%
  \y{Phương trình $L(x)=0$ có hai nghiệm $x=7$ và $x=12$.}{d}
  \y{Với $x=6$ sản phẩm thì bị lỗ.}{d}
  \y{Với $x=9$ sản phẩm thì có lãi.}{d}
  \y{Có lãi với mọi $x\ge7$.}{s}}

\cauDS[id=T10-C6-B17-D73-NB-020,muc=NB,loigiai={Doanh nghiệp không lỗ khi $L(x)\ge0 \Leftrightarrow 2x^2-18x+28\le0$. Tam thức có hai nghiệm $x=2$ và $x=7$, hệ số $a>0$ nên nghiệm của bất phương trình là $2\le x\le7$.}]{Một hợp tác xã trồng rau ở Củ Chi có lợi nhuận khi canh tác $x$ sào đất là $L(x)=-2x^2+18x-28$ (triệu đồng).}{%
  \y{Phương trình $L(x)=0$ có hai nghiệm $x=2$ và $x=7$.}{d}
  \y{Với $x=1$ sào thì bị lỗ.}{d}
  \y{Với $x=4$ sào thì có lãi.}{d}
  \y{Có lãi với mọi $x\ge2$.}{s}}

\cauDS[id=T10-C6-B17-D73-NB-021,muc=NB,loigiai={Doanh nghiệp không lỗ khi $L(x)\ge0 \Leftrightarrow 2x^2-56x+150\le0$. Tam thức có hai nghiệm $x=3$ và $x=25$, hệ số $a>0$ nên nghiệm của bất phương trình là $3\le x\le25$.}]{Một xưởng cơ khí có lợi nhuận khi nhận $x$ đơn hàng mỗi tháng là $L(x)=-2x^2+56x-150$ (triệu đồng).}{%
  \y{Phương trình $L(x)=0$ có hai nghiệm $x=3$ và $x=25$.}{d}
  \y{Với $x=2$ đơn hàng thì bị lỗ.}{d}
  \y{Với $x=14$ đơn hàng thì có lãi.}{d}
  \y{Có lãi với mọi $x\ge3$.}{s}}

\cauDS[id=T10-C6-B17-D73-NB-022,muc=NB,loigiai={Doanh nghiệp không lỗ khi $L(x)\ge0 \Leftrightarrow 2x^2-36x+154\le0$. Tam thức có hai nghiệm $x=7$ và $x=11$, hệ số $a>0$ nên nghiệm của bất phương trình là $7\le x\le11$.}]{Một xưởng cơ khí có lợi nhuận khi nhận $x$ đơn hàng mỗi tháng là $L(x)=-2x^2+36x-154$ (triệu đồng).}{%
  \y{Phương trình $L(x)=0$ có hai nghiệm $x=7$ và $x=11$.}{d}
  \y{Với $x=6$ đơn hàng thì bị lỗ.}{d}
  \y{Với $x=9$ đơn hàng thì có lãi.}{d}
  \y{Có lãi với mọi $x\ge7$.}{s}}

\cauDS[id=T10-C6-B17-D73-NB-023,muc=NB,loigiai={Doanh nghiệp không lỗ khi $L(x)\ge0 \Leftrightarrow x^2-11x+24\le0$. Tam thức có hai nghiệm $x=3$ và $x=8$, hệ số $a>0$ nên nghiệm của bất phương trình là $3\le x\le8$.}]{Một hợp tác xã trồng rau ở Củ Chi có lợi nhuận khi canh tác $x$ sào đất là $L(x)=-x^2+11x-24$ (triệu đồng).}{%
  \y{Phương trình $L(x)=0$ có hai nghiệm $x=3$ và $x=8$.}{d}
  \y{Với $x=2$ sào thì bị lỗ.}{d}
  \y{Với $x=5$ sào thì có lãi.}{d}
  \y{Có lãi với mọi $x\ge3$.}{s}}

\cauDS[id=T10-C6-B17-D73-NB-024,muc=NB,loigiai={Doanh nghiệp không lỗ khi $L(x)\ge0 \Leftrightarrow x^2-46x+304\le0$. Tam thức có hai nghiệm $x=8$ và $x=38$, hệ số $a>0$ nên nghiệm của bất phương trình là $8\le x\le38$.}]{Một quán ăn có lợi nhuận trong ngày khi phục vụ $x$ lượt khách là $L(x)=-x^2+46x-304$ (nghìn đồng).}{%
  \y{Phương trình $L(x)=0$ có hai nghiệm $x=8$ và $x=38$.}{d}
  \y{Với $x=7$ lượt khách thì bị lỗ.}{d}
  \y{Với $x=23$ lượt khách thì có lãi.}{d}
  \y{Có lãi với mọi $x\ge8$.}{s}}

\cauTLN[id=T10-C6-B17-D73-NB-025,muc=NB,loigiai={Doanh nghiệp không lỗ khi $L(x)\ge0 \Leftrightarrow x^2-33x+260\le0$. Tam thức có hai nghiệm $x=13$ và $x=20$, hệ số $a>0$ nên nghiệm của bất phương trình là $13\le x\le20$.}]{Một xưởng cơ khí có lợi nhuận khi nhận $x$ đơn hàng mỗi tháng là $L(x)=-x^2+33x-260$ (triệu đồng). Hỏi số đơn hàng lớn nhất mà vẫn không bị lỗ là bao nhiêu?}{20}

\cauTLN[id=T10-C6-B17-D73-NB-026,muc=NB,loigiai={Doanh nghiệp không lỗ khi $L(x)\ge0 \Leftrightarrow x^2-49x+418\le0$. Tam thức có hai nghiệm $x=11$ và $x=38$, hệ số $a>0$ nên nghiệm của bất phương trình là $11\le x\le38$.}]{Một quán ăn có lợi nhuận trong ngày khi phục vụ $x$ lượt khách là $L(x)=-x^2+49x-418$ (nghìn đồng). Hỏi số lượt khách lớn nhất mà vẫn không bị lỗ là bao nhiêu?}{38}

\cauTLN[id=T10-C6-B17-D73-NB-027,muc=NB,loigiai={Doanh nghiệp không lỗ khi $L(x)\ge0 \Leftrightarrow 2x^2-82x+828\le0$. Tam thức có hai nghiệm $x=18$ và $x=23$, hệ số $a>0$ nên nghiệm của bất phương trình là $18\le x\le23$.}]{Một doanh nghiệp có lợi nhuận khi sản xuất $x$ sản phẩm là $L(x)=-2x^2+82x-828$ (triệu đồng). Hỏi số sản phẩm lớn nhất mà vẫn không bị lỗ là bao nhiêu?}{23}

\cauTLN[id=T10-C6-B17-D73-NB-028,muc=NB,loigiai={Doanh nghiệp không lỗ khi $L(x)\ge0 \Leftrightarrow 2x^2-50x+272\le0$. Tam thức có hai nghiệm $x=8$ và $x=17$, hệ số $a>0$ nên nghiệm của bất phương trình là $8\le x\le17$.}]{Một quán ăn có lợi nhuận trong ngày khi phục vụ $x$ lượt khách là $L(x)=-2x^2+50x-272$ (nghìn đồng). Hỏi số lượt khách lớn nhất mà vẫn không bị lỗ là bao nhiêu?}{17}

\cauTLN[id=T10-C6-B17-D73-NB-029,muc=NB,loigiai={Doanh nghiệp không lỗ khi $L(x)\ge0 \Leftrightarrow x^2-37x+132\le0$. Tam thức có hai nghiệm $x=4$ và $x=33$, hệ số $a>0$ nên nghiệm của bất phương trình là $4\le x\le33$.}]{Một xưởng cơ khí có lợi nhuận khi nhận $x$ đơn hàng mỗi tháng là $L(x)=-x^2+37x-132$ (triệu đồng). Hỏi số đơn hàng lớn nhất mà vẫn không bị lỗ là bao nhiêu?}{33}

\cauTLN[id=T10-C6-B17-D73-NB-030,muc=NB,loigiai={Doanh nghiệp không lỗ khi $L(x)\ge0 \Leftrightarrow 2x^2-52x+176\le0$. Tam thức có hai nghiệm $x=4$ và $x=22$, hệ số $a>0$ nên nghiệm của bất phương trình là $4\le x\le22$.}]{Một quán ăn có lợi nhuận trong ngày khi phục vụ $x$ lượt khách là $L(x)=-2x^2+52x-176$ (nghìn đồng). Hỏi số lượt khách lớn nhất mà vẫn không bị lỗ là bao nhiêu?}{22}

\cauTN[id=T10-C6-B17-D73-TH-001,muc=TH,loigiai={Doanh nghiệp không lỗ khi $L(x)\ge0 \Leftrightarrow 2x^2-60x+400\le0$. Tam thức có hai nghiệm $x=10$ và $x=20$, hệ số $a>0$ nên nghiệm của bất phương trình là $10\le x\le20$.}]{Một xưởng cơ khí có lợi nhuận khi nhận $x$ đơn hàng mỗi tháng là $L(x)=-2x^2+60x-400$ (triệu đồng). Hỏi xưởng cần nhận bao nhiêu đơn hàng để không bị lỗ?}
  {$10\le x\le20$.}
  {$x\le10$ hoặc $x\ge20$.}{$10<x<22$.}{$x\ge20$.}

\cauTN[id=T10-C6-B17-D73-TH-002,muc=TH,loigiai={Doanh nghiệp không lỗ khi $L(x)\ge0 \Leftrightarrow x^2-15x+26\le0$. Tam thức có hai nghiệm $x=2$ và $x=13$, hệ số $a>0$ nên nghiệm của bất phương trình là $2\le x\le13$.}]{Một hợp tác xã trồng rau ở Củ Chi có lợi nhuận khi canh tác $x$ sào đất là $L(x)=-x^2+15x-26$ (triệu đồng). Hỏi hợp tác xã nên canh tác bao nhiêu sào để không bị lỗ?}
  {$2\le x\le13$.}
  {$x\le2$ hoặc $x\ge13$.}{$2<x<15$.}{$x\ge13$.}

\cauTN[id=T10-C6-B17-D73-TH-003,muc=TH,loigiai={Doanh nghiệp không lỗ khi $L(x)\ge0 \Leftrightarrow x^2-63x+962\le0$. Tam thức có hai nghiệm $x=26$ và $x=37$, hệ số $a>0$ nên nghiệm của bất phương trình là $26\le x\le37$.}]{Một xưởng cơ khí có lợi nhuận khi nhận $x$ đơn hàng mỗi tháng là $L(x)=-x^2+63x-962$ (triệu đồng). Hỏi xưởng cần nhận bao nhiêu đơn hàng để không bị lỗ?}
  {$26\le x\le37$.}
  {$x\le26$ hoặc $x\ge37$.}{$26<x<39$.}{$x\ge37$.}

\cauTN[id=T10-C6-B17-D73-TH-004,muc=TH,loigiai={Doanh nghiệp không lỗ khi $L(x)\ge0 \Leftrightarrow 2x^2-66x+432\le0$. Tam thức có hai nghiệm $x=9$ và $x=24$, hệ số $a>0$ nên nghiệm của bất phương trình là $9\le x\le24$.}]{Một quán ăn có lợi nhuận trong ngày khi phục vụ $x$ lượt khách là $L(x)=-2x^2+66x-432$ (nghìn đồng). Hỏi quán cần phục vụ bao nhiêu lượt khách để không bị lỗ?}
  {$9\le x\le24$.}
  {$x\le9$ hoặc $x\ge24$.}{$9<x<26$.}{$x\ge24$.}

\cauTN[id=T10-C6-B17-D73-TH-005,muc=TH,loigiai={Doanh nghiệp không lỗ khi $L(x)\ge0 \Leftrightarrow x^2-46x+513\le0$. Tam thức có hai nghiệm $x=19$ và $x=27$, hệ số $a>0$ nên nghiệm của bất phương trình là $19\le x\le27$.}]{Một hợp tác xã trồng rau ở Củ Chi có lợi nhuận khi canh tác $x$ sào đất là $L(x)=-x^2+46x-513$ (triệu đồng). Hỏi hợp tác xã nên canh tác bao nhiêu sào để không bị lỗ?}
  {$19\le x\le27$.}
  {$x\le19$ hoặc $x\ge27$.}{$19<x<29$.}{$x\ge27$.}

\cauTN[id=T10-C6-B17-D73-TH-006,muc=TH,loigiai={Doanh nghiệp không lỗ khi $L(x)\ge0 \Leftrightarrow 2x^2-106x+1292\le0$. Tam thức có hai nghiệm $x=19$ và $x=34$, hệ số $a>0$ nên nghiệm của bất phương trình là $19\le x\le34$.}]{Một quán ăn có lợi nhuận trong ngày khi phục vụ $x$ lượt khách là $L(x)=-2x^2+106x-1292$ (nghìn đồng). Hỏi quán cần phục vụ bao nhiêu lượt khách để không bị lỗ?}
  {$19\le x\le34$.}
  {$x\le19$ hoặc $x\ge34$.}{$19<x<36$.}{$x\ge34$.}

\cauTN[id=T10-C6-B17-D73-TH-007,muc=TH,loigiai={Doanh nghiệp không lỗ khi $L(x)\ge0 \Leftrightarrow 2x^2-36x+112\le0$. Tam thức có hai nghiệm $x=4$ và $x=14$, hệ số $a>0$ nên nghiệm của bất phương trình là $4\le x\le14$.}]{Một xưởng cơ khí có lợi nhuận khi nhận $x$ đơn hàng mỗi tháng là $L(x)=-2x^2+36x-112$ (triệu đồng). Hỏi xưởng cần nhận bao nhiêu đơn hàng để không bị lỗ?}
  {$4\le x\le14$.}
  {$x\le4$ hoặc $x\ge14$.}{$4<x<16$.}{$x\ge14$.}

\cauTN[id=T10-C6-B17-D73-TH-008,muc=TH,loigiai={Doanh nghiệp không lỗ khi $L(x)\ge0 \Leftrightarrow x^2-57x+800\le0$. Tam thức có hai nghiệm $x=25$ và $x=32$, hệ số $a>0$ nên nghiệm của bất phương trình là $25\le x\le32$.}]{Một quán ăn có lợi nhuận trong ngày khi phục vụ $x$ lượt khách là $L(x)=-x^2+57x-800$ (nghìn đồng). Hỏi quán cần phục vụ bao nhiêu lượt khách để không bị lỗ?}
  {$25\le x\le32$.}
  {$x\le25$ hoặc $x\ge32$.}{$25<x<34$.}{$x\ge32$.}

\cauTN[id=T10-C6-B17-D73-TH-009,muc=TH,loigiai={Doanh nghiệp không lỗ khi $L(x)\ge0 \Leftrightarrow x^2-52x+660\le0$. Tam thức có hai nghiệm $x=22$ và $x=30$, hệ số $a>0$ nên nghiệm của bất phương trình là $22\le x\le30$.}]{Một xưởng cơ khí có lợi nhuận khi nhận $x$ đơn hàng mỗi tháng là $L(x)=-x^2+52x-660$ (triệu đồng). Hỏi xưởng cần nhận bao nhiêu đơn hàng để không bị lỗ?}
  {$22\le x\le30$.}
  {$x\le22$ hoặc $x\ge30$.}{$22<x<32$.}{$x\ge30$.}

\cauTN[id=T10-C6-B17-D73-TH-010,muc=TH,loigiai={Doanh nghiệp không lỗ khi $L(x)\ge0 \Leftrightarrow 2x^2-82x+620\le0$. Tam thức có hai nghiệm $x=10$ và $x=31$, hệ số $a>0$ nên nghiệm của bất phương trình là $10\le x\le31$.}]{Một hợp tác xã trồng rau ở Củ Chi có lợi nhuận khi canh tác $x$ sào đất là $L(x)=-2x^2+82x-620$ (triệu đồng). Hỏi hợp tác xã nên canh tác bao nhiêu sào để không bị lỗ?}
  {$10\le x\le31$.}
  {$x\le10$ hoặc $x\ge31$.}{$10<x<33$.}{$x\ge31$.}

\cauTN[id=T10-C6-B17-D73-TH-011,muc=TH,loigiai={Doanh nghiệp không lỗ khi $L(x)\ge0 \Leftrightarrow 2x^2-46x+252\le0$. Tam thức có hai nghiệm $x=9$ và $x=14$, hệ số $a>0$ nên nghiệm của bất phương trình là $9\le x\le14$.}]{Một quán ăn có lợi nhuận trong ngày khi phục vụ $x$ lượt khách là $L(x)=-2x^2+46x-252$ (nghìn đồng). Hỏi quán cần phục vụ bao nhiêu lượt khách để không bị lỗ?}
  {$9\le x\le14$.}
  {$x\le9$ hoặc $x\ge14$.}{$9<x<16$.}{$x\ge14$.}

\cauTN[id=T10-C6-B17-D73-TH-012,muc=TH,loigiai={Doanh nghiệp không lỗ khi $L(x)\ge0 \Leftrightarrow 2x^2-66x+460\le0$. Tam thức có hai nghiệm $x=10$ và $x=23$, hệ số $a>0$ nên nghiệm của bất phương trình là $10\le x\le23$.}]{Một hợp tác xã trồng rau ở Củ Chi có lợi nhuận khi canh tác $x$ sào đất là $L(x)=-2x^2+66x-460$ (triệu đồng). Hỏi hợp tác xã nên canh tác bao nhiêu sào để không bị lỗ?}
  {$10\le x\le23$.}
  {$x\le10$ hoặc $x\ge23$.}{$10<x<25$.}{$x\ge23$.}

\cauTN[id=T10-C6-B17-D73-TH-013,muc=TH,loigiai={Doanh nghiệp không lỗ khi $L(x)\ge0 \Leftrightarrow x^2-55x+624\le0$. Tam thức có hai nghiệm $x=16$ và $x=39$, hệ số $a>0$ nên nghiệm của bất phương trình là $16\le x\le39$.}]{Một quán ăn có lợi nhuận trong ngày khi phục vụ $x$ lượt khách là $L(x)=-x^2+55x-624$ (nghìn đồng). Hỏi quán cần phục vụ bao nhiêu lượt khách để không bị lỗ?}
  {$16\le x\le39$.}
  {$x\le16$ hoặc $x\ge39$.}{$16<x<41$.}{$x\ge39$.}

\cauTN[id=T10-C6-B17-D73-TH-014,muc=TH,loigiai={Doanh nghiệp không lỗ khi $L(x)\ge0 \Leftrightarrow 2x^2-84x+754\le0$. Tam thức có hai nghiệm $x=13$ và $x=29$, hệ số $a>0$ nên nghiệm của bất phương trình là $13\le x\le29$.}]{Một xưởng cơ khí có lợi nhuận khi nhận $x$ đơn hàng mỗi tháng là $L(x)=-2x^2+84x-754$ (triệu đồng). Hỏi xưởng cần nhận bao nhiêu đơn hàng để không bị lỗ?}
  {$13\le x\le29$.}
  {$x\le13$ hoặc $x\ge29$.}{$13<x<31$.}{$x\ge29$.}

\cauTN[id=T10-C6-B17-D73-TH-015,muc=TH,loigiai={Doanh nghiệp không lỗ khi $L(x)\ge0 \Leftrightarrow 2x^2-122x+1848\le0$. Tam thức có hai nghiệm $x=28$ và $x=33$, hệ số $a>0$ nên nghiệm của bất phương trình là $28\le x\le33$.}]{Một quán ăn có lợi nhuận trong ngày khi phục vụ $x$ lượt khách là $L(x)=-2x^2+122x-1848$ (nghìn đồng). Hỏi quán cần phục vụ bao nhiêu lượt khách để không bị lỗ?}
  {$28\le x\le33$.}
  {$x\le28$ hoặc $x\ge33$.}{$28<x<35$.}{$x\ge33$.}

\cauTN[id=T10-C6-B17-D73-TH-016,muc=TH,loigiai={Doanh nghiệp không lỗ khi $L(x)\ge0 \Leftrightarrow x^2-11x+28\le0$. Tam thức có hai nghiệm $x=4$ và $x=7$, hệ số $a>0$ nên nghiệm của bất phương trình là $4\le x\le7$.}]{Một hợp tác xã trồng rau ở Củ Chi có lợi nhuận khi canh tác $x$ sào đất là $L(x)=-x^2+11x-28$ (triệu đồng). Hỏi hợp tác xã nên canh tác bao nhiêu sào để không bị lỗ?}
  {$4\le x\le7$.}
  {$x\le4$ hoặc $x\ge7$.}{$4<x<9$.}{$x\ge7$.}

\cauTN[id=T10-C6-B17-D73-TH-017,muc=TH,loigiai={Doanh nghiệp không lỗ khi $L(x)\ge0 \Leftrightarrow 2x^2-78x+648\le0$. Tam thức có hai nghiệm $x=12$ và $x=27$, hệ số $a>0$ nên nghiệm của bất phương trình là $12\le x\le27$.}]{Một xưởng cơ khí có lợi nhuận khi nhận $x$ đơn hàng mỗi tháng là $L(x)=-2x^2+78x-648$ (triệu đồng). Hỏi xưởng cần nhận bao nhiêu đơn hàng để không bị lỗ?}
  {$12\le x\le27$.}
  {$x\le12$ hoặc $x\ge27$.}{$12<x<29$.}{$x\ge27$.}

\cauTN[id=T10-C6-B17-D73-TH-018,muc=TH,loigiai={Doanh nghiệp không lỗ khi $L(x)\ge0 \Leftrightarrow x^2-22x+85\le0$. Tam thức có hai nghiệm $x=5$ và $x=17$, hệ số $a>0$ nên nghiệm của bất phương trình là $5\le x\le17$.}]{Một xưởng cơ khí có lợi nhuận khi nhận $x$ đơn hàng mỗi tháng là $L(x)=-x^2+22x-85$ (triệu đồng). Hỏi xưởng cần nhận bao nhiêu đơn hàng để không bị lỗ?}
  {$5\le x\le17$.}
  {$x\le5$ hoặc $x\ge17$.}{$5<x<19$.}{$x\ge17$.}

\cauTN[id=T10-C6-B17-D73-TH-019,muc=TH,loigiai={Doanh nghiệp không lỗ khi $L(x)\ge0 \Leftrightarrow 2x^2-70x+500\le0$. Tam thức có hai nghiệm $x=10$ và $x=25$, hệ số $a>0$ nên nghiệm của bất phương trình là $10\le x\le25$.}]{Một doanh nghiệp có lợi nhuận khi sản xuất $x$ sản phẩm là $L(x)=-2x^2+70x-500$ (triệu đồng). Hỏi doanh nghiệp cần sản xuất bao nhiêu sản phẩm để không bị lỗ?}
  {$10\le x\le25$.}
  {$x\le10$ hoặc $x\ge25$.}{$10<x<27$.}{$x\ge25$.}

\cauTN[id=T10-C6-B17-D73-TH-020,muc=TH,loigiai={Doanh nghiệp không lỗ khi $L(x)\ge0 \Leftrightarrow x^2-77x+1482\le0$. Tam thức có hai nghiệm $x=38$ và $x=39$, hệ số $a>0$ nên nghiệm của bất phương trình là $38\le x\le39$.}]{Một doanh nghiệp có lợi nhuận khi sản xuất $x$ sản phẩm là $L(x)=-x^2+77x-1482$ (triệu đồng). Hỏi doanh nghiệp cần sản xuất bao nhiêu sản phẩm để không bị lỗ?}
  {$38\le x\le39$.}
  {$x\le38$ hoặc $x\ge39$.}{$38<x<41$.}{$x\ge39$.}

\cauTN[id=T10-C6-B17-D73-TH-021,muc=TH,loigiai={Doanh nghiệp không lỗ khi $L(x)\ge0 \Leftrightarrow x^2-54x+585\le0$. Tam thức có hai nghiệm $x=15$ và $x=39$, hệ số $a>0$ nên nghiệm của bất phương trình là $15\le x\le39$.}]{Một quán ăn có lợi nhuận trong ngày khi phục vụ $x$ lượt khách là $L(x)=-x^2+54x-585$ (nghìn đồng). Hỏi quán cần phục vụ bao nhiêu lượt khách để không bị lỗ?}
  {$15\le x\le39$.}
  {$x\le15$ hoặc $x\ge39$.}{$15<x<41$.}{$x\ge39$.}

\cauTN[id=T10-C6-B17-D73-TH-022,muc=TH,loigiai={Doanh nghiệp không lỗ khi $L(x)\ge0 \Leftrightarrow x^2-58x+760\le0$. Tam thức có hai nghiệm $x=20$ và $x=38$, hệ số $a>0$ nên nghiệm của bất phương trình là $20\le x\le38$.}]{Một xưởng cơ khí có lợi nhuận khi nhận $x$ đơn hàng mỗi tháng là $L(x)=-x^2+58x-760$ (triệu đồng). Hỏi xưởng cần nhận bao nhiêu đơn hàng để không bị lỗ?}
  {$20\le x\le38$.}
  {$x\le20$ hoặc $x\ge38$.}{$20<x<40$.}{$x\ge38$.}

\cauTN[id=T10-C6-B17-D73-TH-023,muc=TH,loigiai={Doanh nghiệp không lỗ khi $L(x)\ge0 \Leftrightarrow 2x^2-64x+494\le0$. Tam thức có hai nghiệm $x=13$ và $x=19$, hệ số $a>0$ nên nghiệm của bất phương trình là $13\le x\le19$.}]{Một quán ăn có lợi nhuận trong ngày khi phục vụ $x$ lượt khách là $L(x)=-2x^2+64x-494$ (nghìn đồng). Hỏi quán cần phục vụ bao nhiêu lượt khách để không bị lỗ?}
  {$13\le x\le19$.}
  {$x\le13$ hoặc $x\ge19$.}{$13<x<21$.}{$x\ge19$.}

\cauTN[id=T10-C6-B17-D73-TH-024,muc=TH,loigiai={Doanh nghiệp không lỗ khi $L(x)\ge0 \Leftrightarrow 2x^2-60x+448\le0$. Tam thức có hai nghiệm $x=14$ và $x=16$, hệ số $a>0$ nên nghiệm của bất phương trình là $14\le x\le16$.}]{Một doanh nghiệp có lợi nhuận khi sản xuất $x$ sản phẩm là $L(x)=-2x^2+60x-448$ (triệu đồng). Hỏi doanh nghiệp cần sản xuất bao nhiêu sản phẩm để không bị lỗ?}
  {$14\le x\le16$.}
  {$x\le14$ hoặc $x\ge16$.}{$14<x<18$.}{$x\ge16$.}

\cauDS[id=T10-C6-B17-D73-TH-025,muc=TH,loigiai={Doanh nghiệp không lỗ khi $L(x)\ge0 \Leftrightarrow 2x^2-72x+310\le0$. Tam thức có hai nghiệm $x=5$ và $x=31$, hệ số $a>0$ nên nghiệm của bất phương trình là $5\le x\le31$.}]{Một doanh nghiệp có lợi nhuận khi sản xuất $x$ sản phẩm là $L(x)=-2x^2+72x-310$ (triệu đồng).}{%
  \y{Phương trình $L(x)=0$ có hai nghiệm $x=5$ và $x=31$.}{d}
  \y{Với $x=4$ sản phẩm thì bị lỗ.}{d}
  \y{Với $x=18$ sản phẩm thì có lãi.}{d}
  \y{Có lãi với mọi $x\ge5$.}{s}}

\cauDS[id=T10-C6-B17-D73-TH-026,muc=TH,loigiai={Doanh nghiệp không lỗ khi $L(x)\ge0 \Leftrightarrow 2x^2-70x+132\le0$. Tam thức có hai nghiệm $x=2$ và $x=33$, hệ số $a>0$ nên nghiệm của bất phương trình là $2\le x\le33$.}]{Một hợp tác xã trồng rau ở Củ Chi có lợi nhuận khi canh tác $x$ sào đất là $L(x)=-2x^2+70x-132$ (triệu đồng).}{%
  \y{Phương trình $L(x)=0$ có hai nghiệm $x=2$ và $x=33$.}{d}
  \y{Với $x=1$ sào thì bị lỗ.}{d}
  \y{Với $x=17$ sào thì có lãi.}{d}
  \y{Có lãi với mọi $x\ge2$.}{s}}

\cauDS[id=T10-C6-B17-D73-TH-027,muc=TH,loigiai={Doanh nghiệp không lỗ khi $L(x)\ge0 \Leftrightarrow 2x^2-84x+490\le0$. Tam thức có hai nghiệm $x=7$ và $x=35$, hệ số $a>0$ nên nghiệm của bất phương trình là $7\le x\le35$.}]{Một doanh nghiệp có lợi nhuận khi sản xuất $x$ sản phẩm là $L(x)=-2x^2+84x-490$ (triệu đồng).}{%
  \y{Phương trình $L(x)=0$ có hai nghiệm $x=7$ và $x=35$.}{d}
  \y{Với $x=6$ sản phẩm thì bị lỗ.}{d}
  \y{Với $x=21$ sản phẩm thì có lãi.}{d}
  \y{Có lãi với mọi $x\ge7$.}{s}}

\cauDS[id=T10-C6-B17-D73-TH-028,muc=TH,loigiai={Doanh nghiệp không lỗ khi $L(x)\ge0 \Leftrightarrow 2x^2-78x+496\le0$. Tam thức có hai nghiệm $x=8$ và $x=31$, hệ số $a>0$ nên nghiệm của bất phương trình là $8\le x\le31$.}]{Một xưởng cơ khí có lợi nhuận khi nhận $x$ đơn hàng mỗi tháng là $L(x)=-2x^2+78x-496$ (triệu đồng).}{%
  \y{Phương trình $L(x)=0$ có hai nghiệm $x=8$ và $x=31$.}{d}
  \y{Với $x=7$ đơn hàng thì bị lỗ.}{d}
  \y{Với $x=19$ đơn hàng thì có lãi.}{d}
  \y{Có lãi với mọi $x\ge8$.}{s}}

\cauDS[id=T10-C6-B17-D73-TH-029,muc=TH,loigiai={Doanh nghiệp không lỗ khi $L(x)\ge0 \Leftrightarrow 2x^2-62x+300\le0$. Tam thức có hai nghiệm $x=6$ và $x=25$, hệ số $a>0$ nên nghiệm của bất phương trình là $6\le x\le25$.}]{Một hợp tác xã trồng rau ở Củ Chi có lợi nhuận khi canh tác $x$ sào đất là $L(x)=-2x^2+62x-300$ (triệu đồng).}{%
  \y{Phương trình $L(x)=0$ có hai nghiệm $x=6$ và $x=25$.}{d}
  \y{Với $x=5$ sào thì bị lỗ.}{d}
  \y{Với $x=15$ sào thì có lãi.}{d}
  \y{Có lãi với mọi $x\ge6$.}{s}}

\cauDS[id=T10-C6-B17-D73-TH-030,muc=TH,loigiai={Doanh nghiệp không lỗ khi $L(x)\ge0 \Leftrightarrow 2x^2-48x+256\le0$. Tam thức có hai nghiệm $x=8$ và $x=16$, hệ số $a>0$ nên nghiệm của bất phương trình là $8\le x\le16$.}]{Một xưởng cơ khí có lợi nhuận khi nhận $x$ đơn hàng mỗi tháng là $L(x)=-2x^2+48x-256$ (triệu đồng).}{%
  \y{Phương trình $L(x)=0$ có hai nghiệm $x=8$ và $x=16$.}{d}
  \y{Với $x=7$ đơn hàng thì bị lỗ.}{d}
  \y{Với $x=12$ đơn hàng thì có lãi.}{d}
  \y{Có lãi với mọi $x\ge8$.}{s}}

\cauDS[id=T10-C6-B17-D73-TH-031,muc=TH,loigiai={Doanh nghiệp không lỗ khi $L(x)\ge0 \Leftrightarrow 2x^2-102x+1240\le0$. Tam thức có hai nghiệm $x=20$ và $x=31$, hệ số $a>0$ nên nghiệm của bất phương trình là $20\le x\le31$.}]{Một quán ăn có lợi nhuận trong ngày khi phục vụ $x$ lượt khách là $L(x)=-2x^2+102x-1240$ (nghìn đồng).}{%
  \y{Phương trình $L(x)=0$ có hai nghiệm $x=20$ và $x=31$.}{d}
  \y{Với $x=19$ lượt khách thì bị lỗ.}{d}
  \y{Với $x=25$ lượt khách thì có lãi.}{d}
  \y{Có lãi với mọi $x\ge20$.}{s}}

\cauDS[id=T10-C6-B17-D73-TH-032,muc=TH,loigiai={Doanh nghiệp không lỗ khi $L(x)\ge0 \Leftrightarrow x^2-50x+481\le0$. Tam thức có hai nghiệm $x=13$ và $x=37$, hệ số $a>0$ nên nghiệm của bất phương trình là $13\le x\le37$.}]{Một xưởng cơ khí có lợi nhuận khi nhận $x$ đơn hàng mỗi tháng là $L(x)=-x^2+50x-481$ (triệu đồng).}{%
  \y{Phương trình $L(x)=0$ có hai nghiệm $x=13$ và $x=37$.}{d}
  \y{Với $x=12$ đơn hàng thì bị lỗ.}{d}
  \y{Với $x=25$ đơn hàng thì có lãi.}{d}
  \y{Có lãi với mọi $x\ge13$.}{s}}

\cauTLN[id=T10-C6-B17-D73-TH-033,muc=TH,loigiai={Doanh nghiệp không lỗ khi $L(x)\ge0 \Leftrightarrow x^2-45x+324\le0$. Tam thức có hai nghiệm $x=9$ và $x=36$, hệ số $a>0$ nên nghiệm của bất phương trình là $9\le x\le36$.}]{Một hợp tác xã trồng rau ở Củ Chi có lợi nhuận khi canh tác $x$ sào đất là $L(x)=-x^2+45x-324$ (triệu đồng). Hỏi số sào lớn nhất mà vẫn không bị lỗ là bao nhiêu?}{36}

\cauTLN[id=T10-C6-B17-D73-TH-034,muc=TH,loigiai={Doanh nghiệp không lỗ khi $L(x)\ge0 \Leftrightarrow 2x^2-72x+136\le0$. Tam thức có hai nghiệm $x=2$ và $x=34$, hệ số $a>0$ nên nghiệm của bất phương trình là $2\le x\le34$.}]{Một xưởng cơ khí có lợi nhuận khi nhận $x$ đơn hàng mỗi tháng là $L(x)=-2x^2+72x-136$ (triệu đồng). Hỏi số đơn hàng lớn nhất mà vẫn không bị lỗ là bao nhiêu?}{34}

\cauTLN[id=T10-C6-B17-D73-TH-035,muc=TH,loigiai={Doanh nghiệp không lỗ khi $L(x)\ge0 \Leftrightarrow 2x^2-42x+220\le0$. Tam thức có hai nghiệm $x=10$ và $x=11$, hệ số $a>0$ nên nghiệm của bất phương trình là $10\le x\le11$.}]{Một doanh nghiệp có lợi nhuận khi sản xuất $x$ sản phẩm là $L(x)=-2x^2+42x-220$ (triệu đồng). Hỏi số sản phẩm lớn nhất mà vẫn không bị lỗ là bao nhiêu?}{11}

\cauTLN[id=T10-C6-B17-D73-TH-036,muc=TH,loigiai={Doanh nghiệp không lỗ khi $L(x)\ge0 \Leftrightarrow 2x^2-74x+624\le0$. Tam thức có hai nghiệm $x=13$ và $x=24$, hệ số $a>0$ nên nghiệm của bất phương trình là $13\le x\le24$.}]{Một doanh nghiệp có lợi nhuận khi sản xuất $x$ sản phẩm là $L(x)=-2x^2+74x-624$ (triệu đồng). Hỏi số sản phẩm lớn nhất mà vẫn không bị lỗ là bao nhiêu?}{24}

\cauTLN[id=T10-C6-B17-D73-TH-037,muc=TH,loigiai={Doanh nghiệp không lỗ khi $L(x)\ge0 \Leftrightarrow 2x^2-116x+1482\le0$. Tam thức có hai nghiệm $x=19$ và $x=39$, hệ số $a>0$ nên nghiệm của bất phương trình là $19\le x\le39$.}]{Một doanh nghiệp có lợi nhuận khi sản xuất $x$ sản phẩm là $L(x)=-2x^2+116x-1482$ (triệu đồng). Hỏi số sản phẩm lớn nhất mà vẫn không bị lỗ là bao nhiêu?}{39}

\cauTLN[id=T10-C6-B17-D73-TH-038,muc=TH,loigiai={Doanh nghiệp không lỗ khi $L(x)\ge0 \Leftrightarrow 2x^2-94x+1100\le0$. Tam thức có hai nghiệm $x=22$ và $x=25$, hệ số $a>0$ nên nghiệm của bất phương trình là $22\le x\le25$.}]{Một quán ăn có lợi nhuận trong ngày khi phục vụ $x$ lượt khách là $L(x)=-2x^2+94x-1100$ (nghìn đồng). Hỏi số lượt khách lớn nhất mà vẫn không bị lỗ là bao nhiêu?}{25}

\cauTLN[id=T10-C6-B17-D73-TH-039,muc=TH,loigiai={Doanh nghiệp không lỗ khi $L(x)\ge0 \Leftrightarrow 2x^2-112x+1536\le0$. Tam thức có hai nghiệm $x=24$ và $x=32$, hệ số $a>0$ nên nghiệm của bất phương trình là $24\le x\le32$.}]{Một doanh nghiệp có lợi nhuận khi sản xuất $x$ sản phẩm là $L(x)=-2x^2+112x-1536$ (triệu đồng). Hỏi số sản phẩm lớn nhất mà vẫn không bị lỗ là bao nhiêu?}{32}

\cauTLN[id=T10-C6-B17-D73-TH-040,muc=TH,loigiai={Doanh nghiệp không lỗ khi $L(x)\ge0 \Leftrightarrow 2x^2-80x+672\le0$. Tam thức có hai nghiệm $x=12$ và $x=28$, hệ số $a>0$ nên nghiệm của bất phương trình là $12\le x\le28$.}]{Một xưởng cơ khí có lợi nhuận khi nhận $x$ đơn hàng mỗi tháng là $L(x)=-2x^2+80x-672$ (triệu đồng). Hỏi số đơn hàng lớn nhất mà vẫn không bị lỗ là bao nhiêu?}{28}

\cauTN[id=T10-C6-B17-D73-VD-001,muc=VD,loigiai={Doanh nghiệp không lỗ khi $L(x)\ge0 \Leftrightarrow 2x^2-78x+616\le0$. Tam thức có hai nghiệm $x=11$ và $x=28$, hệ số $a>0$ nên nghiệm của bất phương trình là $11\le x\le28$.}]{Một hợp tác xã trồng rau ở Củ Chi có lợi nhuận khi canh tác $x$ sào đất là $L(x)=-2x^2+78x-616$ (triệu đồng). Hỏi hợp tác xã nên canh tác bao nhiêu sào để không bị lỗ?}
  {$11\le x\le28$.}
  {$x\le11$ hoặc $x\ge28$.}{$11<x<30$.}{$x\ge28$.}

\cauTN[id=T10-C6-B17-D73-VD-002,muc=VD,loigiai={Doanh nghiệp không lỗ khi $L(x)\ge0 \Leftrightarrow x^2-27x+176\le0$. Tam thức có hai nghiệm $x=11$ và $x=16$, hệ số $a>0$ nên nghiệm của bất phương trình là $11\le x\le16$.}]{Một hợp tác xã trồng rau ở Củ Chi có lợi nhuận khi canh tác $x$ sào đất là $L(x)=-x^2+27x-176$ (triệu đồng). Hỏi hợp tác xã nên canh tác bao nhiêu sào để không bị lỗ?}
  {$11\le x\le16$.}
  {$x\le11$ hoặc $x\ge16$.}{$11<x<18$.}{$x\ge16$.}

\cauTN[id=T10-C6-B17-D73-VD-003,muc=VD,loigiai={Doanh nghiệp không lỗ khi $L(x)\ge0 \Leftrightarrow 2x^2-82x+620\le0$. Tam thức có hai nghiệm $x=10$ và $x=31$, hệ số $a>0$ nên nghiệm của bất phương trình là $10\le x\le31$.}]{Một xưởng cơ khí có lợi nhuận khi nhận $x$ đơn hàng mỗi tháng là $L(x)=-2x^2+82x-620$ (triệu đồng). Hỏi xưởng cần nhận bao nhiêu đơn hàng để không bị lỗ?}
  {$10\le x\le31$.}
  {$x\le10$ hoặc $x\ge31$.}{$10<x<33$.}{$x\ge31$.}

\cauTN[id=T10-C6-B17-D73-VD-004,muc=VD,loigiai={Doanh nghiệp không lỗ khi $L(x)\ge0 \Leftrightarrow 2x^2-52x+266\le0$. Tam thức có hai nghiệm $x=7$ và $x=19$, hệ số $a>0$ nên nghiệm của bất phương trình là $7\le x\le19$.}]{Một doanh nghiệp có lợi nhuận khi sản xuất $x$ sản phẩm là $L(x)=-2x^2+52x-266$ (triệu đồng). Hỏi doanh nghiệp cần sản xuất bao nhiêu sản phẩm để không bị lỗ?}
  {$7\le x\le19$.}
  {$x\le7$ hoặc $x\ge19$.}{$7<x<21$.}{$x\ge19$.}

\cauTN[id=T10-C6-B17-D73-VD-005,muc=VD,loigiai={Doanh nghiệp không lỗ khi $L(x)\ge0 \Leftrightarrow x^2-31x+228\le0$. Tam thức có hai nghiệm $x=12$ và $x=19$, hệ số $a>0$ nên nghiệm của bất phương trình là $12\le x\le19$.}]{Một xưởng cơ khí có lợi nhuận khi nhận $x$ đơn hàng mỗi tháng là $L(x)=-x^2+31x-228$ (triệu đồng). Hỏi xưởng cần nhận bao nhiêu đơn hàng để không bị lỗ?}
  {$12\le x\le19$.}
  {$x\le12$ hoặc $x\ge19$.}{$12<x<21$.}{$x\ge19$.}

\cauTN[id=T10-C6-B17-D73-VD-006,muc=VD,loigiai={Doanh nghiệp không lỗ khi $L(x)\ge0 \Leftrightarrow 2x^2-16x+30\le0$. Tam thức có hai nghiệm $x=3$ và $x=5$, hệ số $a>0$ nên nghiệm của bất phương trình là $3\le x\le5$.}]{Một quán ăn có lợi nhuận trong ngày khi phục vụ $x$ lượt khách là $L(x)=-2x^2+16x-30$ (nghìn đồng). Hỏi quán cần phục vụ bao nhiêu lượt khách để không bị lỗ?}
  {$3\le x\le5$.}
  {$x\le3$ hoặc $x\ge5$.}{$3<x<7$.}{$x\ge5$.}

\cauTN[id=T10-C6-B17-D73-VD-007,muc=VD,loigiai={Doanh nghiệp không lỗ khi $L(x)\ge0 \Leftrightarrow 2x^2-96x+1120\le0$. Tam thức có hai nghiệm $x=20$ và $x=28$, hệ số $a>0$ nên nghiệm của bất phương trình là $20\le x\le28$.}]{Một xưởng cơ khí có lợi nhuận khi nhận $x$ đơn hàng mỗi tháng là $L(x)=-2x^2+96x-1120$ (triệu đồng). Hỏi xưởng cần nhận bao nhiêu đơn hàng để không bị lỗ?}
  {$20\le x\le28$.}
  {$x\le20$ hoặc $x\ge28$.}{$20<x<30$.}{$x\ge28$.}

\cauTN[id=T10-C6-B17-D73-VD-008,muc=VD,loigiai={Doanh nghiệp không lỗ khi $L(x)\ge0 \Leftrightarrow x^2-9x+14\le0$. Tam thức có hai nghiệm $x=2$ và $x=7$, hệ số $a>0$ nên nghiệm của bất phương trình là $2\le x\le7$.}]{Một hợp tác xã trồng rau ở Củ Chi có lợi nhuận khi canh tác $x$ sào đất là $L(x)=-x^2+9x-14$ (triệu đồng). Hỏi hợp tác xã nên canh tác bao nhiêu sào để không bị lỗ?}
  {$2\le x\le7$.}
  {$x\le2$ hoặc $x\ge7$.}{$2<x<9$.}{$x\ge7$.}

\cauTN[id=T10-C6-B17-D73-VD-009,muc=VD,loigiai={Doanh nghiệp không lỗ khi $L(x)\ge0 \Leftrightarrow x^2-33x+270\le0$. Tam thức có hai nghiệm $x=15$ và $x=18$, hệ số $a>0$ nên nghiệm của bất phương trình là $15\le x\le18$.}]{Một doanh nghiệp có lợi nhuận khi sản xuất $x$ sản phẩm là $L(x)=-x^2+33x-270$ (triệu đồng). Hỏi doanh nghiệp cần sản xuất bao nhiêu sản phẩm để không bị lỗ?}
  {$15\le x\le18$.}
  {$x\le15$ hoặc $x\ge18$.}{$15<x<20$.}{$x\ge18$.}

\cauTN[id=T10-C6-B17-D73-VD-010,muc=VD,loigiai={Doanh nghiệp không lỗ khi $L(x)\ge0 \Leftrightarrow x^2-35x+304\le0$. Tam thức có hai nghiệm $x=16$ và $x=19$, hệ số $a>0$ nên nghiệm của bất phương trình là $16\le x\le19$.}]{Một hợp tác xã trồng rau ở Củ Chi có lợi nhuận khi canh tác $x$ sào đất là $L(x)=-x^2+35x-304$ (triệu đồng). Hỏi hợp tác xã nên canh tác bao nhiêu sào để không bị lỗ?}
  {$16\le x\le19$.}
  {$x\le16$ hoặc $x\ge19$.}{$16<x<21$.}{$x\ge19$.}

\cauTN[id=T10-C6-B17-D73-VD-011,muc=VD,loigiai={Doanh nghiệp không lỗ khi $L(x)\ge0 \Leftrightarrow 2x^2-146x+2652\le0$. Tam thức có hai nghiệm $x=34$ và $x=39$, hệ số $a>0$ nên nghiệm của bất phương trình là $34\le x\le39$.}]{Một xưởng cơ khí có lợi nhuận khi nhận $x$ đơn hàng mỗi tháng là $L(x)=-2x^2+146x-2652$ (triệu đồng). Hỏi xưởng cần nhận bao nhiêu đơn hàng để không bị lỗ?}
  {$34\le x\le39$.}
  {$x\le34$ hoặc $x\ge39$.}{$34<x<41$.}{$x\ge39$.}

\cauTN[id=T10-C6-B17-D73-VD-012,muc=VD,loigiai={Doanh nghiệp không lỗ khi $L(x)\ge0 \Leftrightarrow x^2-40x+399\le0$. Tam thức có hai nghiệm $x=19$ và $x=21$, hệ số $a>0$ nên nghiệm của bất phương trình là $19\le x\le21$.}]{Một xưởng cơ khí có lợi nhuận khi nhận $x$ đơn hàng mỗi tháng là $L(x)=-x^2+40x-399$ (triệu đồng). Hỏi xưởng cần nhận bao nhiêu đơn hàng để không bị lỗ?}
  {$19\le x\le21$.}
  {$x\le19$ hoặc $x\ge21$.}{$19<x<23$.}{$x\ge21$.}

\cauDS[id=T10-C6-B17-D73-VD-013,muc=VD,loigiai={Doanh nghiệp không lỗ khi $L(x)\ge0 \Leftrightarrow x^2-51x+630\le0$. Tam thức có hai nghiệm $x=21$ và $x=30$, hệ số $a>0$ nên nghiệm của bất phương trình là $21\le x\le30$.}]{Một quán ăn có lợi nhuận trong ngày khi phục vụ $x$ lượt khách là $L(x)=-x^2+51x-630$ (nghìn đồng).}{%
  \y{Phương trình $L(x)=0$ có hai nghiệm $x=21$ và $x=30$.}{d}
  \y{Với $x=20$ lượt khách thì bị lỗ.}{d}
  \y{Với $x=25$ lượt khách thì có lãi.}{d}
  \y{Có lãi với mọi $x\ge21$.}{s}}

\cauDS[id=T10-C6-B17-D73-VD-014,muc=VD,loigiai={Doanh nghiệp không lỗ khi $L(x)\ge0 \Leftrightarrow x^2-46x+408\le0$. Tam thức có hai nghiệm $x=12$ và $x=34$, hệ số $a>0$ nên nghiệm của bất phương trình là $12\le x\le34$.}]{Một doanh nghiệp có lợi nhuận khi sản xuất $x$ sản phẩm là $L(x)=-x^2+46x-408$ (triệu đồng).}{%
  \y{Phương trình $L(x)=0$ có hai nghiệm $x=12$ và $x=34$.}{d}
  \y{Với $x=11$ sản phẩm thì bị lỗ.}{d}
  \y{Với $x=23$ sản phẩm thì có lãi.}{d}
  \y{Có lãi với mọi $x\ge12$.}{s}}

\cauDS[id=T10-C6-B17-D73-VD-015,muc=VD,loigiai={Doanh nghiệp không lỗ khi $L(x)\ge0 \Leftrightarrow x^2-41x+264\le0$. Tam thức có hai nghiệm $x=8$ và $x=33$, hệ số $a>0$ nên nghiệm của bất phương trình là $8\le x\le33$.}]{Một quán ăn có lợi nhuận trong ngày khi phục vụ $x$ lượt khách là $L(x)=-x^2+41x-264$ (nghìn đồng).}{%
  \y{Phương trình $L(x)=0$ có hai nghiệm $x=8$ và $x=33$.}{d}
  \y{Với $x=7$ lượt khách thì bị lỗ.}{d}
  \y{Với $x=20$ lượt khách thì có lãi.}{d}
  \y{Có lãi với mọi $x\ge8$.}{s}}

\cauDS[id=T10-C6-B17-D73-VD-016,muc=VD,loigiai={Doanh nghiệp không lỗ khi $L(x)\ge0 \Leftrightarrow 2x^2-104x+1190\le0$. Tam thức có hai nghiệm $x=17$ và $x=35$, hệ số $a>0$ nên nghiệm của bất phương trình là $17\le x\le35$.}]{Một quán ăn có lợi nhuận trong ngày khi phục vụ $x$ lượt khách là $L(x)=-2x^2+104x-1190$ (nghìn đồng).}{%
  \y{Phương trình $L(x)=0$ có hai nghiệm $x=17$ và $x=35$.}{d}
  \y{Với $x=16$ lượt khách thì bị lỗ.}{d}
  \y{Với $x=26$ lượt khách thì có lãi.}{d}
  \y{Có lãi với mọi $x\ge17$.}{s}}

\cauTLN[id=T10-C6-B17-D73-VD-017,muc=VD,loigiai={Doanh nghiệp không lỗ khi $L(x)\ge0 \Leftrightarrow 2x^2-90x+532\le0$. Tam thức có hai nghiệm $x=7$ và $x=38$, hệ số $a>0$ nên nghiệm của bất phương trình là $7\le x\le38$.}]{Một doanh nghiệp có lợi nhuận khi sản xuất $x$ sản phẩm là $L(x)=-2x^2+90x-532$ (triệu đồng). Hỏi số sản phẩm lớn nhất mà vẫn không bị lỗ là bao nhiêu?}{38}

\cauTLN[id=T10-C6-B17-D73-VD-018,muc=VD,loigiai={Doanh nghiệp không lỗ khi $L(x)\ge0 \Leftrightarrow x^2-31x+168\le0$. Tam thức có hai nghiệm $x=7$ và $x=24$, hệ số $a>0$ nên nghiệm của bất phương trình là $7\le x\le24$.}]{Một xưởng cơ khí có lợi nhuận khi nhận $x$ đơn hàng mỗi tháng là $L(x)=-x^2+31x-168$ (triệu đồng). Hỏi số đơn hàng lớn nhất mà vẫn không bị lỗ là bao nhiêu?}{24}

\cauTLN[id=T10-C6-B17-D73-VD-019,muc=VD,loigiai={Doanh nghiệp không lỗ khi $L(x)\ge0 \Leftrightarrow 2x^2-62x+368\le0$. Tam thức có hai nghiệm $x=8$ và $x=23$, hệ số $a>0$ nên nghiệm của bất phương trình là $8\le x\le23$.}]{Một quán ăn có lợi nhuận trong ngày khi phục vụ $x$ lượt khách là $L(x)=-2x^2+62x-368$ (nghìn đồng). Hỏi số lượt khách lớn nhất mà vẫn không bị lỗ là bao nhiêu?}{23}

\cauTLN[id=T10-C6-B17-D73-VD-020,muc=VD,loigiai={Doanh nghiệp không lỗ khi $L(x)\ge0 \Leftrightarrow 2x^2-76x+210\le0$. Tam thức có hai nghiệm $x=3$ và $x=35$, hệ số $a>0$ nên nghiệm của bất phương trình là $3\le x\le35$.}]{Một xưởng cơ khí có lợi nhuận khi nhận $x$ đơn hàng mỗi tháng là $L(x)=-2x^2+76x-210$ (triệu đồng). Hỏi số đơn hàng lớn nhất mà vẫn không bị lỗ là bao nhiêu?}{35}

\cauTN[id=T10-C6-B17-D73-VDC-001,muc=VDC,loigiai={Doanh nghiệp không lỗ khi $L(x)\ge0 \Leftrightarrow 2x^2-116x+1482\le0$. Tam thức có hai nghiệm $x=19$ và $x=39$, hệ số $a>0$ nên nghiệm của bất phương trình là $19\le x\le39$.}]{Một xưởng cơ khí có lợi nhuận khi nhận $x$ đơn hàng mỗi tháng là $L(x)=-2x^2+116x-1482$ (triệu đồng). Hỏi xưởng cần nhận bao nhiêu đơn hàng để không bị lỗ?}
  {$19\le x\le39$.}
  {$x\le19$ hoặc $x\ge39$.}{$19<x<41$.}{$x\ge39$.}

\cauTN[id=T10-C6-B17-D73-VDC-002,muc=VDC,loigiai={Doanh nghiệp không lỗ khi $L(x)\ge0 \Leftrightarrow 2x^2-68x+290\le0$. Tam thức có hai nghiệm $x=5$ và $x=29$, hệ số $a>0$ nên nghiệm của bất phương trình là $5\le x\le29$.}]{Một xưởng cơ khí có lợi nhuận khi nhận $x$ đơn hàng mỗi tháng là $L(x)=-2x^2+68x-290$ (triệu đồng). Hỏi xưởng cần nhận bao nhiêu đơn hàng để không bị lỗ?}
  {$5\le x\le29$.}
  {$x\le5$ hoặc $x\ge29$.}{$5<x<31$.}{$x\ge29$.}

\cauTN[id=T10-C6-B17-D73-VDC-003,muc=VDC,loigiai={Doanh nghiệp không lỗ khi $L(x)\ge0 \Leftrightarrow 2x^2-40x+192\le0$. Tam thức có hai nghiệm $x=8$ và $x=12$, hệ số $a>0$ nên nghiệm của bất phương trình là $8\le x\le12$.}]{Một quán ăn có lợi nhuận trong ngày khi phục vụ $x$ lượt khách là $L(x)=-2x^2+40x-192$ (nghìn đồng). Hỏi quán cần phục vụ bao nhiêu lượt khách để không bị lỗ?}
  {$8\le x\le12$.}
  {$x\le8$ hoặc $x\ge12$.}{$8<x<14$.}{$x\ge12$.}

\cauTN[id=T10-C6-B17-D73-VDC-004,muc=VDC,loigiai={Doanh nghiệp không lỗ khi $L(x)\ge0 \Leftrightarrow x^2-17x+70\le0$. Tam thức có hai nghiệm $x=7$ và $x=10$, hệ số $a>0$ nên nghiệm của bất phương trình là $7\le x\le10$.}]{Một doanh nghiệp có lợi nhuận khi sản xuất $x$ sản phẩm là $L(x)=-x^2+17x-70$ (triệu đồng). Hỏi doanh nghiệp cần sản xuất bao nhiêu sản phẩm để không bị lỗ?}
  {$7\le x\le10$.}
  {$x\le7$ hoặc $x\ge10$.}{$7<x<12$.}{$x\ge10$.}

\cauTN[id=T10-C6-B17-D73-VDC-005,muc=VDC,loigiai={Doanh nghiệp không lỗ khi $L(x)\ge0 \Leftrightarrow x^2-60x+851\le0$. Tam thức có hai nghiệm $x=23$ và $x=37$, hệ số $a>0$ nên nghiệm của bất phương trình là $23\le x\le37$.}]{Một doanh nghiệp có lợi nhuận khi sản xuất $x$ sản phẩm là $L(x)=-x^2+60x-851$ (triệu đồng). Hỏi doanh nghiệp cần sản xuất bao nhiêu sản phẩm để không bị lỗ?}
  {$23\le x\le37$.}
  {$x\le23$ hoặc $x\ge37$.}{$23<x<39$.}{$x\ge37$.}

\cauTN[id=T10-C6-B17-D73-VDC-006,muc=VDC,loigiai={Doanh nghiệp không lỗ khi $L(x)\ge0 \Leftrightarrow x^2-45x+506\le0$. Tam thức có hai nghiệm $x=22$ và $x=23$, hệ số $a>0$ nên nghiệm của bất phương trình là $22\le x\le23$.}]{Một quán ăn có lợi nhuận trong ngày khi phục vụ $x$ lượt khách là $L(x)=-x^2+45x-506$ (nghìn đồng). Hỏi quán cần phục vụ bao nhiêu lượt khách để không bị lỗ?}
  {$22\le x\le23$.}
  {$x\le22$ hoặc $x\ge23$.}{$22<x<25$.}{$x\ge23$.}

\cauDS[id=T10-C6-B17-D73-VDC-007,muc=VDC,loigiai={Doanh nghiệp không lỗ khi $L(x)\ge0 \Leftrightarrow 2x^2-50x+272\le0$. Tam thức có hai nghiệm $x=8$ và $x=17$, hệ số $a>0$ nên nghiệm của bất phương trình là $8\le x\le17$.}]{Một doanh nghiệp có lợi nhuận khi sản xuất $x$ sản phẩm là $L(x)=-2x^2+50x-272$ (triệu đồng).}{%
  \y{Phương trình $L(x)=0$ có hai nghiệm $x=8$ và $x=17$.}{d}
  \y{Với $x=7$ sản phẩm thì bị lỗ.}{d}
  \y{Với $x=12$ sản phẩm thì có lãi.}{d}
  \y{Có lãi với mọi $x\ge8$.}{s}}

\cauDS[id=T10-C6-B17-D73-VDC-008,muc=VDC,loigiai={Doanh nghiệp không lỗ khi $L(x)\ge0 \Leftrightarrow x^2-39x+140\le0$. Tam thức có hai nghiệm $x=4$ và $x=35$, hệ số $a>0$ nên nghiệm của bất phương trình là $4\le x\le35$.}]{Một quán ăn có lợi nhuận trong ngày khi phục vụ $x$ lượt khách là $L(x)=-x^2+39x-140$ (nghìn đồng).}{%
  \y{Phương trình $L(x)=0$ có hai nghiệm $x=4$ và $x=35$.}{d}
  \y{Với $x=3$ lượt khách thì bị lỗ.}{d}
  \y{Với $x=19$ lượt khách thì có lãi.}{d}
  \y{Có lãi với mọi $x\ge4$.}{s}}

\cauTLN[id=T10-C6-B17-D73-VDC-009,muc=VDC,loigiai={Doanh nghiệp không lỗ khi $L(x)\ge0 \Leftrightarrow 2x^2-70x+500\le0$. Tam thức có hai nghiệm $x=10$ và $x=25$, hệ số $a>0$ nên nghiệm của bất phương trình là $10\le x\le25$.}]{Một quán ăn có lợi nhuận trong ngày khi phục vụ $x$ lượt khách là $L(x)=-2x^2+70x-500$ (nghìn đồng). Hỏi số lượt khách lớn nhất mà vẫn không bị lỗ là bao nhiêu?}{25}

\cauTLN[id=T10-C6-B17-D73-VDC-010,muc=VDC,loigiai={Doanh nghiệp không lỗ khi $L(x)\ge0 \Leftrightarrow x^2-24x+140\le0$. Tam thức có hai nghiệm $x=10$ và $x=14$, hệ số $a>0$ nên nghiệm của bất phương trình là $10\le x\le14$.}]{Một xưởng cơ khí có lợi nhuận khi nhận $x$ đơn hàng mỗi tháng là $L(x)=-x^2+24x-140$ (triệu đồng). Hỏi số đơn hàng lớn nhất mà vẫn không bị lỗ là bao nhiêu?}{14}
